\begin{textsample}{NEG Dim 3 – global\_south\_2024\_12 – Score 1.00 – global\_south\_2024\_12/118345387.txt}  \label{ex:f3_neg_068}
FILE PHOTO : Displaced Palestinians make their way as they flee areas in northern Gaza Strip following an Israeli evacuation order, amid the Israel-Hamas conflict, in Jabalia, October 6, 2024.
Photograph : ( Reuters ) Story highlights We are pausing the delivery of aid through Kerem Shalom, the main crossing point for humanitarian aid into Gaza.
The road out of this crossing has not been safe for months.
On 16 November, a large convoy of aid trucks was stolen by armed gangs, UNRWA head Philippe Lazzarini said in a post on X.
The UN agency for Palestinian refugees has halted the delivery of aid through the key Kerem Shalom crossing between Israel and Gaza after it became"impossible", its chief said Sunday."We are pausing the delivery of aid through Kerem Shalom, the main crossing point for humanitarian aid into Gaza.
The road out of this crossing has not been safe for months.
On 16 November, a large convoy of aid trucks was stolen by armed gangs,"UNRWA head Philippe Lazzarini said in a post on X."Yesterday, we tried to bring in a few food trucks on the same route.
They were all taken,"he added, warning hunger was"rapidly deepening"in Gaza.
Lazzarini listed how the humanitarian operation had become"unnecessarily impossible"due to"the ongoing siege, hurdles from Israeli authorities, political decisions to restrict the amounts of aid, lack of safety on aid routes and targeting of local police".
It comes after the United Nations warned Friday that Gaza has descended into anarchy, with hunger soaring, looting rampant and rising numbers of rapes in shelters as public order falls apart.
Lazzarini called on Israel to ensure aid flowed to Gaza and said the country"must refrain from attacks on humanitarian workers".
According to UNRWA ' s count, only 65 aid trucks per day had been able to enter Gaza the past month, compared to a pre-war average of 500.
Israel, which imposed a total siege on the Hamas-ruled territory in the early stages of the war last year, blames the inability of relief organisations to handle and distribute large quantities of aid. ' Catastrophic ' During a press visit Thursday, the Israeli army showed aid shipments at the crossing and said they wait at the Gaza side of Kerem Shalom for"months".
Lazzarini ' s announcement follows an Israeli strike Saturday that killed three contractors of the US charity World Central Kitchen, including one who Israel ' s military said was involved in Hamas ' s October 7, 2023 attack.
The United Nations said last month that 333 aid workers had been killed since the start of the war in October of last year, 243 of them employees of UNRWA.
Lazzarini reiterated his call for a ceasefire"that would also secure the delivery of safe and uninterrupted aid to people in need".
Jean-Francois Corty, president of France-based charity Medecins du Monde, warned UNRWA ' s decision was a"very bad omen"and"tragic in a context that was already so".
Claire Nicolet, head of mission in Jerusalem for Doctors Without Borders, told AFP the situation was"already catastrophic"and that UNRWA ' s announcement was the"straw that broke the camel ' s back"because the UN agency was"the backbone of aid for the supply of food and equipment".
The Hamas attack on October 7, 2023 resulted in the deaths of 1,207 people on Israeli side, most of them civilians, according to an AFP tally based on official Israeli figures.
At least 44,429 Palestinians, a majority of them civilians, have been killed in Israel ' s military campaign in the Gaza Strip since the war began, according to data provided by the Hamas-run health ministry in the territory.
The UN has acknowledged these figures as reliable.
Disclaimer : This story has been published from a news agency feed with minimal edits to adhere to WION ' s style guide.
The headline may have been changed to better reflect the content of the story or to make it more suitable for WION audience.

% matched lemmas: 
\end{textsample}
